\documentclass[conference]{IEEEtran}
\usepackage{amsmath}
\usepackage{booktabs}
\usepackage{cite}
\author{\IEEEauthorblockN{Ritik, Aman, Ansh Tomar, Sarabjeet Singh, and Deepanshu Chaudhary}}


\title{Trace-Driven Evaluation of a Two-Way LRU Cache\\
\large Phase II Project Report}
\begin{document}
\maketitle

\begin{abstract}
We compare two Verilog caches with equal 64-byte capacity and four-byte
blocks. The baseline is direct mapped. The second design has two ways per
set and uses least recently used (LRU) replacement. Across seven shared
traces, two-way LRU raises the alternating-conflict hit rate from 0\% to
87.5\% and reduces total service cycles from 112 to 42. With a cyclic
three-block trace, direct mapping retains one of the blocks and produces
seven hits, while two-way LRU produces zero hits. Generic Yosys synthesis
maps the direct and two-way controllers to 1,782 and 1,920 cells,
respectively. The results show how the distribution of requests across
sets affects both hit rate and service time.
\end{abstract}

\begin{IEEEkeywords}
Cache memory, Verilog, set associativity, LRU, trace-driven simulation
\end{IEEEkeywords}

\section{Introduction and Related Work}
Under direct mapping, every memory block has one eligible cache line.
Set associativity gives a block more eligible slots within a set. At fixed
data capacity, increasing the number of ways also reduces the number of
sets, which changes which blocks compete for those slots. Smith surveys
these placement and replacement choices alongside block size and write
policy \cite{smith}. Hill and Smith compare alternative cache organizations
using identical address traces to study the effect of associativity
on misses \cite{hill}. We apply that approach to two small hardware
controllers and report both access behavior and the hardware cost of the
implemented designs.

\section{Implementation}
Each controller accepts one CPU request at a time. Both use 16-bit byte
addresses, four-byte blocks, 8-bit CPU data, write-through, and
write-allocate. Direct mapping has 16 lines and uses address bits
$[5:2]$ as its index. Two-way LRU has eight sets and uses $[4:2]$ to
select a set. Both use $[1:0]$ to select a byte. The high bits form the tag.

On each request, the two-way controller checks both ways for a valid tag
match. On a miss, it fills an invalid way if one is available. For a full
set, a stored bit identifies the LRU victim. Each hit or fill sets that
bit to the other way. The controller records the original hit decision
before a refill. On a write miss it installs the fetched block with the
selected byte updated. It then writes that byte through to memory and
waits for acknowledgment. The CPU and memory each use separate request
and completion signals.

\begin{table}[!t]
\caption{Service Cycles at Three-Cycle Memory Delay}
\label{tab:cycles}
\centering
\begin{tabular}{lccc}
\toprule
Access & Block fetches & Byte writes & Cycles\\
\midrule
Read hit & 0 & 0 & 2\\
Read miss & 1 & 0 & 7\\
Write hit & 0 & 1 & 7\\
Write miss & 1 & 1 & 12\\
\bottomrule
\end{tabular}
\end{table}

\section{Experimental Method}
Each of the seven traces starts with empty cache lines and identically
initialized byte-addressable memory. Both caches execute the same
requests in order at memory latency $L=3$. We count hits from the initial
lookup and measure service cycles from request acceptance to response,
excluding idle gaps. For a request that incurs $N_{\mathrm{mem}}$ memory
transactions, the cost is
\[
 C = 2 + (L + 2)N_{\mathrm{mem}}.
\]
The two base cycles are lookup and response. Each memory transaction adds
its delay and two handshake cycles. Table~\ref{tab:cycles} lists the four
access types for the comparison.

The testbench independently tracks the stored bytes and evaluates
replacement using timestamps. It checks cycle counts, memory operations,
and cache results. The RTL passed 3,259 completed requests per cache at
each of three memory delays (1, 3, and 7 cycles). Coverage includes
writes to every byte offset, readbacks, address boundaries, LRU
replacement, and reset during an outstanding refill. A separate Python
model checks returned data, hit decisions, and service cycles for all
902 saved events from the 14 trace runs.

\begin{table}[!t]
\caption{Measured Hit Rates and Total Service Cycles}
\label{tab:results}
\centering
\footnotesize
\setlength{\tabcolsep}{2.2pt}
\begin{tabular}{lrrrrr}
\toprule
& & \multicolumn{2}{c}{Hit rate (\%)} & \multicolumn{2}{c}{Cycles}\\
\cmidrule(lr){3-4}\cmidrule(l){5-6}
Workload & Reqs. & Direct & Two-way & Direct & Two-way\\
\midrule
Alternating conflict & 16 & 0.000 & 87.500 & 112 & 42\\
Repeated address & 16 & 93.750 & 93.750 & 37 & 37\\
Sequential bytes & 128 & 87.500 & 87.500 & 336 & 336\\
LRU replacement & 7 & 0.000 & 42.857 & 49 & 34\\
Capacity overflow & 64 & 0.000 & 0.000 & 448 & 448\\
Three colliding blocks & 24 & 29.167 & 0.000 & 133 & 168\\
Mixed reads and writes & 196 & 64.796 & 60.714 & 1077 & 1117\\
\bottomrule
\end{tabular}
\end{table}

\section{Results and Discussion}
\subsection{Trace Comparisons}
In Table~\ref{tab:results}, the alternating-conflict trace switches
between \texttt{0x0000} and \texttt{0x0040}. Both addresses map to direct
line 0. Two-way LRU keeps the pair in set 0 and hits on all 14 requests
after the first two fills. Total service cycles fall by 62.5\% for
this trace.

In the three-block trace, \texttt{0x0000}, \texttt{0x0020}, and
\texttt{0x0040} repeat eight times. All three blocks map to two-way
set 0. LRU evicts a block before its next access on every iteration.
In direct mapping, \texttt{0x0020} occupies line 8 and hits on its
last seven accesses. The mixed workload also favors direct mapping
by eight hits and 40 cycles. Repeated and sequential workloads tie
because both caches retain the blocks involved in reuse. The capacity
trace cycles twice through 32 blocks with space for only 16.

\begin{table}[!t]
\caption{Storage and Generic Synthesis Counts}
\label{tab:synth}
\centering
\begin{tabular}{lrr}
\toprule
Measure & Direct & Two-way\\
\midrule
Data bits & 512 & 512\\
Tag bits & 160 & 176\\
Valid bits & 16 & 16\\
LRU bits & 0 & 8\\
Generic cells & 1,782 & 1,920\\
\bottomrule
\end{tabular}
\end{table}

\newpage
\subsection{Hardware Cost and Scope}
Both controllers were mapped through the same Yosys generic synthesis
flow. Each generated netlist passed the 3,259-request test at latency 3.
Table~\ref{tab:synth} shows 138 additional generic cells for two-way LRU
(7.7\%). The flow maps storage to registers and logic. The counts
describe that mapping. A target-specific implementation is needed to
measure FPGA utilization or maximum clock frequency.

The workloads are short and synthetic. The design accepts one outstanding
request, and the memory model uses a fixed latency with separate block
reads and byte writes. The cycle counts apply to these controllers
and the stated memory model. The experiment does not model a processor
pipeline or a physical memory bus.


\end{document}
